% ===========================================================================
\section{Bounded differences and martingales}\label{sec:bdd}
\warmth{0}\quad\whisper{Upgrade ``sum of independents'' to ``any function that changes little when one coordinate moves''.}

\begin{strand}
In one line: it need not be a sum. As long as $f(x_1,\dots,x_n)$ changes by a bounded amount
when a single coordinate is altered, reveal its randomness \emph{one coordinate at a time}
(the Doob martingale); each step has bounded difference, feed it to Hoeffding's lemma from
§\ref{sec:subgaussian}, and concentration appears.
\end{strand}

\block{The paradigm}
\trigger{$f$ is not a sum, but has \keyword{bounded differences}: altering coordinate $i$ changes it by $|f-f'|\le c_i$.}
Build the \keyword{Doob martingale} $M_k=\E[f\mid X_1,\dots,X_k]$, so $M_0=\E f$, $M_n=f$, with
increments $|M_k-M_{k-1}|\le c_k$. \pick{subgaussian}Apply Hoeffding's lemma per step and sum.

\begin{strand}
Three-line heuristic: (1) \keyword{any} function is its own Doob martingale (revealing one coordinate = one step);
(2) concentration depends only on \emph{how far each step can jump}, not on how complex $f$ is;
(3) the martingale constant ($\sum c_i^2$) is looser than the entropy method (§\ref{sec:entropy}) by a factor,
but the proof is shortest and the assumptions weakest.
\end{strand}

\block{Cheat-sheet: the bounded-difference family}
\noindent\begin{tabularx}{\linewidth}{@{}l l L@{}}
\toprule
{\color{indigo}Theorem} & {\color{indigo}Setting} & \multicolumn{1}{l}{\color{indigo}Tail bound / when to use}\\
\midrule
Azuma–Hoeffding & increments $|M_k-M_{k-1}|\le c_k$ & $\Pr[|M_n-M_0|\ge t]\le 2e^{-t^2/2\sum c_k^2}$\\
McDiarmid       & $f$ bounded diff.\ $c_i$, $X_i$ indep.\ & $\Pr[|f-\E f|\ge t]\le 2e^{-2t^2/\sum c_i^2}$; workhorse of ML generalization\\
Self-bounding   & $f$ self-bounding              & one-sided Bernstein type, variance-aware (cf.\ §\ref{sec:entropy})\\
Freedman        & martingale + predictable var.\ $V$ & variance-aware martingale tail, tighter when $V$ small\\
\bottomrule
\end{tabularx}

\block{Main results \& proof skeletons}

\begin{theorem}[Azuma–Hoeffding]\label{thm:azuma}
$(M_k)$ a martingale, $|M_k-M_{k-1}|\le c_k$ a.s. Then $\Pr[M_n-M_0\ge t]\le \exp\!\big(-t^2/2\textstyle\sum_k c_k^2\big)$.
\end{theorem}
\begin{proof}
Skeleton: for the increments $D_k=M_k-M_{k-1}$ we have $\E[D_k\mid\mathcal F_{k-1}]=0$ and $|D_k|\le c_k$,
so the conditional Hoeffding lemma gives $\E[e^{\lambda D_k}\mid\mathcal F_{k-1}]\le e^{\lambda^2 c_k^2/2}$;
peel off layer by layer via the tower property to get $\E e^{\lambda(M_n-M_0)}\le e^{\lambda^2\sum c_k^2/2}$; then Chernoff.
\TODO{fill the conditional Hoeffding lemma and the tower-property unrolling}
\end{proof}

\begin{theorem}[McDiarmid's bounded-differences inequality]\label{thm:mcdiarmid}
$X_1,\dots,X_n$ independent, $f$ with bounded differences $c_i$. Then $\Pr[|f-\E f|\ge t]\le 2\exp\!\big(-2t^2/\textstyle\sum_i c_i^2\big)$.
\end{theorem}
\begin{proof}\TODO{feed the Doob martingale $M_k=\E[f\mid X_{1:k}]$ into Thm~\ref{thm:azuma}; independence keeps each increment's range $\le c_k$ (sup minus inf)}\end{proof}

\begin{example}[Concentration of the generalization gap]\label{ex:generalization}
$f=\sup_{h\in\mathcal H}\big(\E_{\mathcal D}\ell_h-\widehat{\E}\,\ell_h\big)$ has bounded differences $\le 1/n$ in each sample,
so McDiarmid gives $\Pr[f-\E f\ge t]\le e^{-2nt^2}$. Thus the generalization gap concentrates uniformly around its
expectation ($=$ the Rademacher complexity). \TODO{verify the bounded-difference constant $1/n$}
\end{example}

\begin{strand}
Why it goes this way: McDiarmid splits ``uniform convergence'' into two things, \emph{concentration}
(this section, $f$ hugs $\E f$) and \emph{the size of the expectation} ($\E f=$ Rademacher complexity, estimated
by the chaining of §\ref{sec:chaining}). That is the source of the standard two-step recipe in statistical learning theory.
\end{strand}

\loose{Bounded differences can be relaxed: the asymmetric / variance-aware refinement (BLM Ch.\,6) replaces $\sum c_i^2$ by the true variance, but needs the entropy method. Room to improve here.}%
\warp{tensorize} The truly tight bound comes from the \emph{tensorization} of §\ref{sec:entropy}; this section is its ``martingale approximation''. \pick{tensorize}
